\phantomsection\addcontentsline{toc}{chapter}{Introduction} % Don't remove me!
\chapter*{Introduction}
The brain is among the most widely studied subjects of research, and yet many neural behaviours remain elusive to scientists. Vision is an important neural process, and provides one of the first gates of attention to the human brain. A large part of the brain is devoted to vision, the visual cortex containing between 4-6 billion of the 30 billion neurons of the neocortex, up to 5 times the size of other primates \cite{wandell_visual_2009, andrews_correlated_1997}. In particular, search tasks reveal complex visual dynamics and invite research into mechanisms of attention and brain function through a functional understanding of gaze patterns. To study these patterns of gaze dynamics, methods adapted from statistical physics, network science, and non-equilibrium thermodynamics have been found to explain and provide understanding of neural computation. The resulting approaches to understanding individual patterns of human gaze during visual search have found, most often through training large neural networks to predict gaze trajectories \cite{bhandari_modeling_2025, wu_gitnet_2026,abdelrahman_deep_2026,bremer_predicting_2021,ozdel_gaze-guided_2024,akinyelu_convolutional_2022} and saliency maps from an image for human observers \cite{kummerer_modeling_2025,cornia_deep_2016, cornia_predicting_2018, kammerer2016deepgaze2, huang_salicon_2015, bujia_modeling_2022}. Such study has fostered developments in the broader field of artificial intelligence, where inspiration from the behaviour of the human brain is often explicit. Most notably, the DeepGaze III project developed by  \textcite{Kummerer2022_DeepGazeIII} has yielded high accuracy in predicting the location of the fixations of human vision given the image viewed, the fixation history, and biased patterns of searching \cite{Kummerer2022_DeepGazeIII, kummerer_modeling_2025, dagostino2025what}. Such DeepGaze models have allowed researchers to understand how simple low-level contrast is often enough to predict the location of fixations \cite{kummerer_understanding_2017}, and how late fixations quickly become dominated by gaze bias and decouple from the geometry of the image viewed \cite{Kummerer2022-au}.

Although, a well-motivated dynamical account of visual search dynamics beyond fixation prediction remains unexplored. Such a unifying formulation should account for a large range of well-documented patterns of visual search behaviour. For instance, search dynamics are identified by the spatiotemporal patterns of the  excursions of visual focus called saccades, the dynamical behaviour during fixation between such attentional saccades, the preferential searching after interesting areas of high saliency, the inhibition of attention in previously explored areas, and the emergence of timescales of attention. To be valuable, such a formulation should be biologically plausible, and yield testable predictions and functional explanations of novel or unexplained aspects of visual dynamics.

% Such approaches give often quite accurate results, but they lack the rich understanding of a full model of the network of the brain itself. Theories of biological attention and cognition, such as those of the ``Bayesian brain'', built on active sensing and inference provide a more motivated framework to understand the brain as maintaining a rational causal mental worldview \cite{knill_bayesian_2004, costa_probabilistic_2024, bujia_modeling_2022, tan_theta_2024}. This Bayesian model of the brain has been a fruitful direction to understand the most basic level of neural cognition as sampling through population activity patterns \cite{orban_neural_2016, NIPS2002_a486cd07, beck_probabilistic_2008, zhu_autocorrelated_2024}. Although it has been a debated framework (see \cite{mangalam_myth_2025}), the Bayesian brain has provided a strong window into biological cognition. Geoffrey Hinton and colleagues invented the Helmholtz machine, performing explicit Bayesian inference on perceptive inputs, and updating predictions on the external world, minimising the free energy as a loss function of a neural network. Such a formulation inspired pioneering work on neural coding \cite{Lange2023-sn, beck_probabilistic_2008}, Karl Friston's highly influential Free Energy Principle in theoretical neuroscience \cite{Friston2010-aa, Parr2025-xh}, and the resulting empirical models and well-verified predictions about the human sensory-motor system \cite{kording_bayesian_2004,rao_predictive_1999,koch_predicting_1999}.